
\documentclass{article}
\usepackage{colortbl}
\usepackage{makecell}
\usepackage{multirow}
\usepackage{supertabular}

\begin{document}

\newcounter{utterance}

\centering \large Interaction Transcript for game `cladder', experiment `full\_v1.5\_default', episode 5310 with Qwen3.5{-}4B.
\vspace{24pt}

{ \footnotesize  \setcounter{utterance}{1}
\setlength{\tabcolsep}{0pt}
\begin{supertabular}{c@{$\;$}|p{.15\linewidth}@{}p{.15\linewidth}p{.15\linewidth}p{.15\linewidth}p{.15\linewidth}p{.15\linewidth}}
    \# & \multicolumn{2}{c}{Player} && \multicolumn{2}{c}{Game Master} \\
    \hline

    \theutterance \stepcounter{utterance}  
    & & & \multicolumn{4}{p{0.6\linewidth}}{
        \cellcolor[rgb]{0.9,0.9,0.9}{
            \makecell[{{p{\linewidth}}}]{
                \texttt{\tiny{[P1$\langle$GM]}}
                \texttt{You are an expert in causal inference. The following question is not a typical commonsense query, but rather a meticulously designed question created by a professor specializing in causal inference, intended to assess the students' mastery of the course content.} \\
\texttt{Imagine a self{-}contained, hypothetical world with only the following conditions, and without any unmentioned factors or causal relationships: Liking spicy food has a direct effect on occupation and salary. Skill has a direct effect on occupation and salary. Occupation has a direct effect on salary. Skill is unobserved. Method 1: We look at how liking spicy food correlates with salary case by case according to occupation. Method 2: We look directly at how liking spicy food correlates with salary in general. To understand how liking spicy food affects salary, is it more correct to use the Method 1 than Method 2?} \\
\texttt{Start your answer with 'Yes' or 'No', followed by additional reasoning or evidence to support your explanation.} \\
            }
        }
    }
     \\ \\

    \theutterance \stepcounter{utterance}  
    & \multicolumn{4}{p{0.6\linewidth}}{
        \cellcolor[rgb]{0.561,0.737,0.561}{
            \makecell[{{p{\linewidth}}}]{
                \texttt{\tiny{[P1$\rangle$GM]}}
                \texttt{Yes} \\
\\ 
\texttt{**Reasoning:**} \\
\\ 
\texttt{To determine the correct method, we must analyze the causal structure described:} \\
\texttt{1.  **Spicy Food (\$S\$)** \$\textbackslash{}rightarrow\$ **Occupation (\$O\$)**} \\
\texttt{2.  **Skill (\$U\$)** \$\textbackslash{}rightarrow\$ **Occupation (\$O\$)**} \\
\texttt{3.  **Skill (\$U\$)** \$\textbackslash{}rightarrow\$ **Salary (\$Y\$)**} \\
\texttt{4.  **Occupation (\$O\$)** \$\textbackslash{}rightarrow\$ **Salary (\$Y\$)**} \\
\\ 
\texttt{This structure creates a **backdoor path** from Spicy Food to Salary:} \\
\texttt{\$S \textbackslash{}rightarrow O \textbackslash{}leftarrow U \textbackslash{}rightarrow Y\$} \\
\\ 
\texttt{Here, **Occupation (\$O\$)** acts as a **confounder** (or more precisely in this specific directed acyclic graph, a collider on the path \$S \textbackslash{}rightarrow O \textbackslash{}leftarrow U\$, but since \$U\$ also affects \$Y\$, the path \$S \textbackslash{}rightarrow O \textbackslash{}leftarrow U \textbackslash{}rightarrow Y\$ is an open backdoor path that induces a spurious correlation between \$S\$ and \$Y\$).} \\
\\ 
\texttt{*   **Method 2 (General Correlation):** This method calculates the total correlation between \$S\$ and \$Y\$ without conditioning. Because \$S\$ and \$Y\$ share a common cause (\$U\$, Skill) indirectly through \$O\$, and because \$S\$ affects \$O\$ which affects \$Y\$, the observed correlation in Method 2 is a mixture of the direct causal effect (\$S \textbackslash{}rightarrow O \textbackslash{}rightarrow Y\$) and the spurious association induced by the backdoor path (\$S \textbackslash{}rightarrow O \textbackslash{}leftarrow U \textbackslash{}rightarrow Y\$). Without controlling for the confounding variable (Occupation, which mediates the effect of \$S\$ but is also influenced by the confounder \$U\$), Method 2 yields a biased estimate of the causal effect.} \\
\\ 
\texttt{*   **Method 1 (Stratification by Occupation):** This method calculates the correlation between \$S\$ and \$Y\$ within each specific level of Occupation (\$O\$). By conditioning on \$O\$, we block the backdoor path (\$S \textbackslash{}rightarrow O \textbackslash{}leftarrow U \textbackslash{}rightarrow Y\$). Within a fixed occupation, the influence of \$U\$ on the relationship between \$S\$ and \$Y\$ is severed (or at least, the variation due to \$O\$ is removed). This allows us to isolate the effect of \$S\$ on \$Y\$ that operates *through* the occupation channel, effectively estimating the causal effect \$P(Y|do(S))\$ (assuming no other unobserved confounders exist between \$S\$ and \$Y\$ directly or via other paths).} \\
\\ 
\texttt{Therefore, because Method 1 blocks the confounding bias introduced by the unobserved Skill variable via the Occupation variable, it provides a more accurate estimate of the causal effect than the unadjusted general correlation in Method 2.} \\
            }
        }
    }
    & & \\ \\

    \theutterance \stepcounter{utterance}  
    & & & \multicolumn{2}{p{0.3\linewidth}}{
        \cellcolor[rgb]{0.9,0.9,0.9}{
            \makecell[{{p{\linewidth}}}]{
                \texttt{\tiny{[GM$|$GM]}}
                \texttt{Yes} \\
            }
        }
    }
    & & \\ \\

    \theutterance \stepcounter{utterance}  
    & & & \multicolumn{2}{p{0.3\linewidth}}{
        \cellcolor[rgb]{0.9,0.9,0.9}{
            \makecell[{{p{\linewidth}}}]{
                \texttt{\tiny{[GM$|$GM]}}
                \texttt{no} \\
            }
        }
    }
    & & \\ \\

    \theutterance \stepcounter{utterance}  
    & & & \multicolumn{2}{p{0.3\linewidth}}{
        \cellcolor[rgb]{0.9,0.9,0.9}{
            \makecell[{{p{\linewidth}}}]{
                \texttt{\tiny{[GM$|$GM]}}
                \texttt{game\_result = LOSE} \\
            }
        }
    }
    & & \\ \\

\end{supertabular}
}

\end{document}
